\documentclass{article}
\usepackage[french]{babel}
\usepackage[utf8]{inputenc}
\usepackage[T1]{fontenc}
\usepackage{amsmath}
\usepackage{amssymb}
\usepackage{graphicx}
\usepackage{float}
\usepackage{hyperref}
\usepackage{subcaption}

%%%%%%%%%%%%%%%% Lengths %%%%%%%%%%%%%%%%
\setlength{\textwidth}{15.5cm}
\setlength{\evensidemargin}{0.5cm}
\setlength{\oddsidemargin}{0.5cm}

%%%%%%%%%%%%%%%% Variables %%%%%%%%%%%%%%%%
\def\projet{6}
\def\titre{Résolution approchée d'équations différentielles / Modélisation de systèmes dynamiques}
\def\groupe{2}
\def\equipe{3}
\def\responsible{Benjamin Chew}
\def\secretary{Solan Boher--Sanchez}
\def\others{Alexis Gaudray Bouju, Yani Haouchine}

\begin{document}

%%%%%%%%%%%%%%%% Header %%%%%%%%%%%%%%%%
\noindent\begin{minipage}{0.98\textwidth}
  \vskip 0mm
  \noindent
  { \begin{tabular}{p{7.5cm}}
      {\bfseries \sffamily
        Projet \projet} \\
      {\itshape \titre}
    \end{tabular}}
  \hfill
  \fbox{\begin{tabular}{l}
      {~\hfill \bfseries \sffamily Groupe \groupe\ - Equipe \equipe \hfill~} \\[2mm]
      Responsable : \responsible \\
      Secrétaire : \secretary \\
      Codeurs : \others
    \end{tabular}}
  \vskip 4mm ~

  ~~~\parbox{0.95\textwidth}{\small \textit{Résumé~: Ce projet a pour objectif la résolution approchée d'équations différentielles àl'aide de différentes méthodes et son application à des systèmes physiques complexes. La première application est l'analyse du comportement chaotique d'un double pendule. La seconde est le système proie-prédateur de Lotka-Volterra, qui est un système modélisant l'évolution d'une population. } \sffamily

  }
  \vskip 1mm ~
\end{minipage}

%%%%%%%%%%%%%%%% Main part %%%%%%%%%%%%%%%%

\section{Résolution d'équations différentielles}
Afin de résoudre une équation différentielle, nous allons implémenter et comparer plusieurs méthodes de Runge-Kutta. Tout d'abord la méthode d'Euler puis la méthode du point milieu, celle de Heun et enfin la méthode de Runge-Kutta d'ordre 4.

\subsection{Présentation des méthodes}
Dans cette section, nous allons détailler les différentes méthodes de Runge-Kutta que nous avons implémentées.

\subsubsection{Méthode d'Euler}
La méthode d'Euler (rectangle à gauche) est la plus simple des méthodes de Runge-Kutta. On évalue la fonction à l'extrémité gauche de l'intervalle en \(x=0\)
La formule de la méthode d'Euler est donnée par :
\(y_{n+1} = y_n + h_n F(t_n, y_n)\)
où \(h_n\) est le pas de temps, \(t_n\) est le temps courant et \(y_n\) est la valeur de la fonction à ce temps.

\subsubsection{Méthode du point milieu}
La méthode du point milieu est une amélioration de la méthode d'Euler. On évalue la fonction à au milieu de l'intervalle en \(x=\frac{1}{2}\)
La formule de la méthode du point milieu est donnée par :
\(y_{n+1} = y_n + h_n F(t_n + \frac{h_n}{2}, y_n + \frac{h_n}{2} F(t_n, y_n))\)

\subsubsection{Méthode de Heun}
La méthode de Heun est l'équivalent de la méthode du trapèze. Elle est d'ordre 2. La formule de la méthode de Heun est donnée par :
\(y_{n+1} = y_n + \frac{h_n}{2} ( F(t_n, y_n) + F(t_n + h_n, y_n + h_n F(t_n, y_n)) )\)

\subsubsection{Méthode de Runge-Kutta d'ordre 4}
La méthode de Runge-Kutta d'ordre 4 est une méthode plus avancée d'ordre 4. Elle correspond à la méthode de Simpson. La formule de la méthode de Runge-Kutta d'ordre 4 est donnée par :
\(y_{n+1} = y_n + \frac{h_n}{6} ( p_{n1} + 2p_{n2} + 2p_{n3} + p_{n4} )\) où :
\(p_{n1} = F(t_n, y_n)\)
\(p_{n2} = F(t_n + \frac{h_n}{2}, y_n + \frac{h_n}{2} p_{n1})\)
\(p_{n3} = F(t_n + \frac{h_n}{2}, y_n + \frac{h_n}{2} p_{n2})\)
\(p_{n4} = F(t_n + h_n, y_n + h_n p_{n3})\)

\subsection{Comparaison des méthodes}
Nous allons comparer les différentes méthodes de Runge-Kutta en traçant la solution de l'équation différentielle pour chaque méthode. Nous allons également calculer l'erreur entre la solution exacte et la solution approchée pour chaque méthode.

Nous allons tester ces fonctions sur l'équation différentielle suivante dont la solution exacte est connue :
\(y(t) = \left[ \begin{array}{c} y_1(t) \\ y_2(t) \end{array} \right]\)
, avec :
\(\begin{cases} y(0) = \left[ \begin{array}{c} 1 \\ 0 \end{array} \right] \\ y’(t) = \left[ \begin{array}{c} -y_2(t) \\ y_1(t) \end{array} \right] \end{cases}\)
La solution est :
\(y(t) = \left[ \begin{array}{c} \cos(t) \\ \sin(t) \end{array} \right] \)

\begin{figure}[H]
  \centering
  \begin{subfigure}[t]{0.55\textwidth}
    \centering
    \includegraphics[width=\textwidth]{rapport/img/champ_tangente.png}
    \caption{Champ des tangentes.}
    \label{fig:champ_tangente_a}
  \end{subfigure}
  \hfill
  \begin{subfigure}[t]{0.43\textwidth}
    \centering
    \includegraphics[width=\textwidth]{rapport/img/comparaison_methods.png}
    \caption{Comparaison des méthodes.}
    \label{fig:comparaison_methods}
  \end{subfigure}
  \caption{Exemple sur une équation à solution connue}
  \label{fig:champs_tangentes_comparaison}
\end{figure}

Nous observons sur le tracé des différentes méthodes de Runge-Kutta Figure~\ref{fig:comparaison_methods} que la méthode d'Euler est la moins bonne et diverge. Les autres méthodes (superposées sur le Figure~\ref{fig:comparaison_methods}) sont plus précises et convergent vers la solution exacte. Ces solutions concordent avec le champ des tangentes de la solution exacte, tracé sur la Figure~\ref{fig:champ_tangente_a}.

\section{Étude du pendule à $N$ maillons}
Nous considérons un système composé de $N$ tiges rigides, de longueur $\ell$ et de masse $m$, reliées en série et libres de se mouvoir dans un plan vertical. Chaque tige peut pivoter autour de l’extrémité de la précédente. Ce système est appelé \textit{pendule à $N$ maillons}, et son état est entièrement décrit par un vecteur d’angles :
\[
\boldsymbol{\theta} = (\theta_1, \theta_2, \dots, \theta_N)
\]
où $\theta_i$ est l’angle formé entre la tige $i$ et la verticale.

\subsection{Pendule simple}
Le pendule simple obéit à l’équation différentielle non linéaire : \(\ddot{\theta} + \frac{g}{\ell} \sin(\theta) = 0\)
\\
Pour de petites oscillations ($\theta \ll 1$), on peut linéariser l’équation : \(\ddot{\theta} + \frac{g}{\ell} \theta = 0\).
\\
Ce qui correspond à un oscillateur harmonique de fréquence : \(\omega = \sqrt{\frac{g}{\ell}}\)

Dans la Figure~\ref{fig:pendulum_frequency} on observe bien que pour des petites oscillations la fréquence est égale à ${\sqrt{\frac{g}{l}}}$.

\begin{figure}[H]
  \centering
  \includegraphics[width=0.57\textwidth]{rapport/img/pendulum_frequency.png}
  \caption{Fréquence du pendule à un unique maillon.}
  \label{fig:pendulum_frequency}
\end{figure}
\subsection{Double pendule}
Le pendule à deux maillons est un système à deux degrés de liberté. Il obéit à un système de deux équations différentielles couplées. Les équations sont non linéaires, fortement couplées et sensibles aux conditions initiales.
La figure \ref{fig:extremity_double_pendulum} représente la trajectoire de l'extrémité du pendule pour une vitesse initiale nulle et des angles $\theta_1, \theta_2$ initiaux non nulles. Le système présente l'aspect chaotique caractéristique du double pendule.
\begin{figure}[H]
  \centering
  \includegraphics[width=0.62\textwidth]{rapport/img/extremity_double_pendulum.png}
  \caption{Trajectoire de l'extrémité du pendule à deux maillons.}
  \label{fig:extremity_double_pendulum}
\end{figure}
\subsubsection*{Sensibilité aux conditions initiales}
Pour illustrer le caractère chaotique on a représenté en figure \ref{fig:first_flip_time} le temps pris par le système pour un premier retournement, c'est à dire $\theta_1$ ou $\theta_2$ supérieure à un angle pi, en fonction de différents angles initiaux.

On remarque que plus les angles initiaux sont simultanément proches de 0 plus le premier retournement prend du temps. Cependant, un écart important entre les 2 angles initiaux ne garantie pas un retournement rapide.
\begin{figure}[H]
  \centering
  \includegraphics[width=0.62\textwidth]{rapport/img/first_flip_time.png}
  \caption{Temps de premier retournement en fonction des angles initiaux.}
  \label{fig:first_flip_time}
\end{figure}

\section*{Simulation d'un modéle proie-predateur}
Il s'agit ici de simuler le systéme proie-prédateur, à l'aide du modéle de Lotka-Volterra :

\begin{equation}
\begin{cases}
\displaystyle \frac{dN}{dt} = aN - bNP \\
\displaystyle \frac{dP}{dt} = cNP - dP
\end{cases}
\end{equation}

où :
\begin{itemize}
\item \( N(t) \) : population de proies au temps \( t \),
\item \( P(t) \) : population de prédateurs au temps \( t \),
\item \( a > 0 \) : taux de croissance des proies en l'absence de prédateurs,
\item \( b > 0 \) : taux de mortalité des proies dû aux prédateurs,
\item \( c > 0 \) : taux de croissance des prédateurs en fonction de la disponibilité des proies,
\item \( d > 0 \) : taux de mortalité naturelle des prédateurs.
\end{itemize}

\subsection*{Variation au cours du temps}

\begin{figure}[H]
  \centering
  \includegraphics[width=0.7\textwidth]{rapport/img/modéle_proie_predateur.png}
  \caption{Évolution au cours du temps de nombres de proie et du nombre de prédateur à droite et portrait de phase à droite}
  \label{fig:population_evolution}
\end{figure}

La figure \ref{fig:population_evolution} à été construite en choisissons a = 1.2 , b = 0.3 , c = 0.8 et d = 1.0
Ces valeurs sont cohérentes avec un systéme ou les proies seraient des lapins et les prédateurs seraient des renards

\end{document}
